% !TEX root = ../main.tex

\section{Triad Construction}

A \textbf{Triad} is a 3-note chord built by stacking thirds (skipping every other letter in the scale).

\begin{itemize}
    \item \textbf{Scale degrees:} 1 (Root) -- 3 (Third) -- 5 (Fifth).
    \item \textbf{Major Triad:} Root $+ 4$ half-steps (Major 3rd) $+ 3$ half-steps (Minor 3rd).
    \item \textbf{Minor Triad:} Root $+ 3$ half-steps (Minor 3rd) $+ 4$ half-steps (Major 3rd).
\end{itemize}

\begin{examplebox}{The Major to Minor Rule}
To turn any Major triad into a Minor triad, \textbf{lower the 3rd (middle note) by 1 half-step} to the left:
\begin{itemize}
    \item C Major ($C - E - G$) $\longrightarrow$ C minor ($C - E\flat - G$)
    \item E Major ($E - G\sharp - B$) $\longrightarrow$ E minor ($E - G - B$)
    \item G Major ($G - B - D$) $\longrightarrow$ G minor ($G - B\flat - D$)
\end{itemize}
\end{examplebox}

\section{Why Notes are Called 3rd, 5th, 7th}

Numbers denote the \textbf{interval distance / degree in the parent scale}, not physical key count or finger order:
\begin{itemize}
    \item \textbf{Root (1st / Tonic):} Foundation note.
    \item \textbf{3rd (Mediant):} Defines whether the chord sounds Major or Minor.
    \item \textbf{5th (Dominant):} Provides harmonic stability and frame.
\end{itemize}

\begin{definitionbox}{Interval Number vs. Physical Keys (Semitones)}
Why is $G\sharp$ the "Third" of $E$ Major when it is physically 4 keys away on the keyboard?
\begin{itemize}
    \item \textbf{The Number counts Alphabet Letters:} Starting from $E$:
    \begin{equation*}
        \text{Letter 1: } E \quad\longrightarrow\quad \text{Letter 2: } F \quad\longrightarrow\quad \text{Letter 3: } G
    \end{equation*}
    Any letter $G$ ($G$, $G\sharp$, $G\flat$) is always a \textbf{Third} above $E$.
    \item \textbf{Physical Semitones define Quality (Major vs. Minor):}
    \begin{itemize}
        \item $E \rightarrow G$: 3 half-steps $\longrightarrow$ \textbf{Minor 3rd} (smaller 3rd).
        \item $E \rightarrow G\sharp$: 4 half-steps $\longrightarrow$ \textbf{Major 3rd} (larger 3rd).
    \end{itemize}
    \item \textbf{Scale Degree:} In the E Major scale ($E - F\sharp - \mathbf{G\sharp} - A - B - C\sharp - D\sharp$), $G\sharp$ is literally the 3rd note.
\end{itemize}
\end{definitionbox}

\section{Triad Inversions}

Inversions change the lowest (bass) note by rolling the bottom note up one octave:

\begin{itemize}
    \item \textbf{Root Position:} Root at bottom ($1 - 3 - 5$). Example: $C - E - G$.
    \item \textbf{1st Inversion:} 3rd at bottom ($3 - 5 - 1$). Example: $E - G - C$.
    \item \textbf{2nd Inversion:} 5th at bottom ($5 - 1 - 3$). Example: $G - C - E$.
\end{itemize}

\begin{warningbox}{Triad Inversion Limits}
A 3-note triad only has 3 possible bottom notes, meaning it only has Root Position, 1st Inversion, and 2nd Inversion. A 3rd Inversion only exists for 4-note chords (7ths).
\end{warningbox}

\section{Other Chord Types \& Notations}

\subsection{Suspended Chords (sus2, sus4)}

The 3rd is omitted and replaced by the 2nd or 4th scale degree (neither major nor minor):
\begin{itemize}
    \item \textbf{sus4:} $1 - 4 - 5$ (Root $+ 5$ semitones $+ 2$ semitones).
    \begin{itemize}
        \item Example: $\text{Csus4} = C - F - G$ (resolves $F \rightarrow E$).
    \end{itemize}
    \item \textbf{sus2:} $1 - 2 - 5$ (Root $+ 2$ semitones $+ 5$ semitones).
    \begin{itemize}
        \item Example: $\text{Csus2} = C - D - G$ (open sound).
    \end{itemize}
\end{itemize}

\subsection{Augmented Triad (aug / +)}

Formed by stacking two Major 3rds:
\begin{itemize}
    \item \textbf{Formula:} $1 - 3 - \sharp 5$ (Root $+ 4$ semitones $+ 4$ semitones).
    \item Example: $C\text{aug} = C - E - G\sharp$
\end{itemize}

\subsection{Added Tone Chords (add9)}

A standard triad with the 9th degree added on top (without including a 7th):
\begin{itemize}
    \item \textbf{Formula:} $1 - 3 - 5 - 9$ (9 is the 2nd note an octave higher).
    \item Example: $C\text{add}9 = C - E - G - D$
    \item Contrast with $C9$: $C9$ is a Dominant 7th with an added 9th ($C - E - G - B\flat - D$).
\end{itemize}

\subsection{Slash Chords (Chord / Bass)}

Notation: $\text{Right-Hand Chord} \;/\; \text{Left-Hand Single Bass Note}$.
\begin{itemize}
    \item \textbf{Inversion Notation:} $C / E$ = C Major triad in RH ($C - E - G$) over an $E$ bass note (1st Inversion).
    \item \textbf{Polychord / Step Bass:} $B\flat / A\flat$ = B$\flat$ Major triad in RH ($B\flat - D - F$) over an $A\flat$ bass note.
\end{itemize}
